% Simulation environment section for OrbitGraph paper.

\section{Simulation Environment}
\label{sec:simulation}

Evaluating routing under realistic LEO dynamics requires a platform that can
(i) evolve topology and link delays as satellites move,
(ii) exercise real intra-domain routing or an equivalent control-plane install
path, and
(iii) exchange packets end-to-end with measurable latency and loss.
We surveyed representative LEO simulators and emulators following a recent
ns-3-based framework review~\cite{manzanareslopez2025comprehensive} and our
project evaluation notes.
Table~\ref{tab:simulators} summarizes the alternatives we considered; we adopt
\textbf{StarryNet}~\cite{lai2023starrynet} for the experiments in this paper.

\begin{table}[t]
  \centering
  \caption{LEO simulation platforms considered for OrbitGraph evaluation.}
  \label{tab:simulators}
  \renewcommand{\arraystretch}{1.15}
  \begin{tabular}{@{}lccc@{}}
    \hline
    Platform & Packet-level & Real routing & Live topology \\
    \hline
    SILLEO-SCNS~\cite{kempton2020silleo}      & No  & No  & Yes \\
    Hypatia~\cite{kassing2020exploring}       & Yes & No  & Yes \\
    LEOPath~\cite{leopath2024}                & No  & No  & Yes \\
    StarPerf~\cite{lai2020starperf}           & No  & No  & Yes \\
    StarryNet~\cite{lai2023starrynet}         & Yes & Yes & Yes \\
    \hline
  \end{tabular}
  \\[0.5em]
  \footnotesize
  ``Real routing'': distributed protocols (e.g., OSPF) or kernel FIB install on
  emulated nodes, not precomputed abstract paths alone.
\end{table}

\subsection{Candidate platforms and selection rationale}

\paragraph{SILLEO-SCNS.}
SILLEO-SCNS (Simulator for Large LEO Satellite Communication
Networks)~\cite{kempton2020silleo} focuses on constellation generation, orbital
motion, and inter-satellite connectivity modeling.
It offers strong topology visualization and scales well for studying link
geometry.
However, it does not perform packet-level network simulation: there is no
end-to-end traffic exchange through a real or emulated protocol stack, so it
cannot directly measure RTT, jitter, or sustained data-plane outage around
routing events---metrics central to our OrbitGraph evaluation.

\paragraph{Hypatia.}
Hypatia~\cite{kassing2020exploring} integrates precomputed orbital state with
the ns-3 discrete-event simulator and supports packet-level experiments at LEO
scale.
Its routing layer, however, relies on precomputed shortest paths (e.g.,
Floyd--Warshall) rather than running distributed protocols such as OSPF or an
operator-style centralized install.
That design is well suited to capacity and latency characterization but does
not let us compare distributed reconvergence against proactive SDN-style
updates on identical emulated topologies.
In addition, full ns-3 campaign runtimes reported in the literature are
substantially longer than container emulation for comparable scenario
complexity.

\paragraph{LEOPath.}
LEOPath~\cite{leopath2024} is a Python framework for analyzing and comparing
routing algorithms over time-varying LEO graphs.
It is more extensible than static shortest-path pipelines and targets routing
research directly.
LEOPath does not integrate ns-3 and does not drive end-to-end packet exchange
through kernel forwarding; routing decisions are evaluated in a custom
simulation core rather than against live TCP/IP stacks and real handover
timing on emulated links.

\paragraph{StarPerf (StarPath).}
StarPerf~\cite{lai2020starperf}, sometimes referred to as StarPath in
community repositories, improves on Hypatia-style workflows with faster
area-to-area performance profiling and constellation scaling.
It excels at statistical characterization of mega-constellation design
trade-offs under abstract traffic models.
Like Hypatia and LEOPath, it does not instantiate real protocol stacks on
emulated nodes; custom routing algorithms are evaluated in a flow- or
performance-oriented model rather than through kernel routes installed before
a ground--satellite link teardown.

\paragraph{StarryNet (chosen).}
StarryNet~\cite{lai2023starrynet} emulates integrated space--terrestrial
networks with Docker containers and Linux traffic control.
Each satellite and ground station runs in its own container; inter-satellite
and ground--satellite links are Docker bridge networks whose propagation delay,
loss, and bandwidth are enforced with \texttt{tc netem} in the Linux kernel.
Orbital motion is driven by SGP4-based ephemerides; per-second delay matrices
and a \texttt{Topo\_leo\_change} schedule capture continuous topology drift
and GSL handovers.
Crucially for our study, StarryNet supports \emph{real} intra-domain routing:
the baseline runs BIRD with OSPF, while our OrbitGraph controller installs
kernel static routes via \texttt{ip route replace}---the same data plane a
centralized SDN deployment would program, without modifying application code.
End-to-end probes (\texttt{ping}, continuous outage sampling) traverse the
actual forwarding path, so control-plane actions are reflected in measurable
data-plane availability.

We therefore use StarryNet as the evaluation harness: it is the only surveyed
platform that combines time-varying LEO topology, live link mutation at
handover ticks, and comparable OSPF versus centralized route-install
experiments on identical emulated graphs.

\subsection{Emulation model used in our experiments}

Within StarryNet we configure Walker-style $O \times S$ grids with plus-grid
inter-satellite links, two fixed ground stations, LeastDelay link selection,
550\,km altitude, and 53$^\circ$ inclination---representative of published
Starlink-class parameters.
Each node is a privileged Linux container (\texttt{lwsen/starlab\_node:1.0});
ISLs and GSLs appear and disappear according to the precomputed topology
change log, while delay matrices refresh every 10\,s to reflect orbital motion.
The emulation loop executes in wall-clock time: routing events (damage,
recovery, GSL handover) are applied synchronously at scheduled seconds, matching
StarryNet's intended mode for protocol and application testing.

Our OrbitGraph extensions (incremental FIB deltas, make-before-break install
ordering, and proactive handover) sit entirely in the control plane; the
underlying orbital, link, and container model remains StarryNet's published
design~\cite{lai2023starrynet}.
